% !TeX program = xelatex
\documentclass[12pt,a4paper]{article}
\usepackage{fontspec}
\usepackage{polyglossia}
\setmainlanguage{russian}
\setmainfont{DejaVu Serif}
\setmonofont{DejaVu Sans Mono}
\usepackage[a4paper,margin=2cm]{geometry}
\usepackage{microtype}
\usepackage{enumitem}
\usepackage{hyperref}
\usepackage{booktabs}
\usepackage{tabularx}
\usepackage{longtable}
\usepackage{listings}
\hypersetup{colorlinks=true,linkcolor=blue!50!black,urlcolor=blue!60!black}
\lstset{basicstyle=\ttfamily\small,breaklines=true,frame=single,columns=fullflexible,keepspaces=true,showstringspaces=false,tabsize=4}
\title{Конспект первой лекции\\\large «Проектирование архитектуры веб-приложений»}
\date{}
\begin{document}
\maketitle
\tableofcontents
\newpage

\section{Общие сведения}
\subsection{Тема}
Варианты архитектур веб-приложений. От простого серверного приложения к клиент--серверным, модульным и микросервисным системам. Мастер-класс создания простейшего REST API на FastAPI с ORM, SQLite, валидацией и пагинацией.

\subsection{Продолжительность}
3 академических часа.

\subsection{Связь с практической работой №1}
Практическая работа №1 требует создать REST API-сервис с CRUD, SQLite или другой СУБД, ORM, полной валидацией и частичной загрузкой данных. Поэтому мастер-класс строится как полностью готовый пример именно такого сервиса.

На первой лекции рассматриваются \textbf{только веб-приложения}. React не вводится: он появляется на следующем этапе, когда уже имеется готовый API.

\section{План занятия}
\begin{tabularx}{\textwidth}{@{}lXr@{}}
\toprule
Время & Содержание & Длительность\\
\midrule
0:00--0:10 & Что такое архитектура веб-приложения & 10 мин\\
0:10--0:55 & Основные варианты архитектур веб-приложений & 45 мин\\
0:55--1:10 & Сравнение и выбор архитектуры курса & 15 мин\\
1:10--1:25 & HTTP, REST и CRUD & 15 мин\\
1:25--1:35 & Создание проекта и первый endpoint & 10 мин\\
1:35--1:50 & Pydantic и первый API & 15 мин\\
1:50--2:10 & SQLite и SQLAlchemy ORM & 20 мин\\
2:10--2:35 & CRUD & 25 мин\\
2:35--2:55 & Пагинация, валидация и проверка & 20 мин\\
2:55--3:00 & Связь с практической работой №1 & 5 мин\\
\bottomrule
\end{tabularx}

\section{Часть I. Варианты архитектур веб-приложений}
\subsection{Что такое архитектура}
Архитектура отвечает на вопрос: как части веб-приложения разделены и как они взаимодействуют.

Для веб-приложения обычно рассматриваются браузер или другой HTTP-клиент, серверная часть, база данных и внешние сервисы. Фреймворк не является архитектурой сам по себе: FastAPI, Django или Spring являются инструментами реализации.

\subsection{Сервер генерирует HTML}
Простейшая схема:
\begin{verbatim}
Браузер --HTTP--> Web-приложение --SQL--> База данных
                    |
                    +---- HTML ----> Браузер
\end{verbatim}
Сервер получает запрос, выполняет логику, обращается к БД и формирует HTML-страницу.

\subsection{Многослойное серверное приложение}
Даже один серверный процесс можно логически разделить:
\begin{verbatim}
HTTP
 |
v
Presentation
 |
v
Application / Business Logic
 |
v
Data Access
 |
v
Database
\end{verbatim}
Важно: логическое разделение слоёв не означает обязательного разделения на отдельные процессы или серверы.

\subsection{SPA + REST API}
\begin{verbatim}
Браузер / SPA --HTTP+JSON--> REST API --ORM--> База данных
\end{verbatim}
Сервер отдаёт данные, а интерфейс формируется клиентским приложением. API потенциально может использоваться разными клиентами.

\subsection{Монолит}
Монолит -- единое разворачиваемое серверное приложение. Оно может содержать API, бизнес-логику и доступ к данным внутри одного процесса.

\begin{verbatim}
+----------------------------------+
|          Web application         |
| API | Business | Data Access     |
+----------------+-----------------+
                 |
                 v
              Database
\end{verbatim}

Монолит не является синонимом плохой архитектуры. Для небольшой системы он часто является самым простым решением.

\subsection{Модульный монолит}
Функциональность разделена на модули, но приложение по-прежнему разворачивается как единое целое:
\begin{verbatim}
+---------------------------------------+
|              Application             |
|  +---------+ +---------+ +---------+ |
|  | Users   | | Orders  | | Catalog | |
|  +---------+ +---------+ +---------+ |
+-------------------+-------------------+
                    |
                    v
                 Database
\end{verbatim}

\subsection{Микросервисная архитектура}
Функциональность разделяется на самостоятельно разворачиваемые сервисы:
\begin{verbatim}
                  Client
                    |
                    v
               API Gateway
              /     |      \\
             v      v       v
          Users   Orders   Catalog
            |        |        |
           DB       DB       DB
\end{verbatim}
Распределённость даёт независимость развёртывания, но создаёт сетевые ошибки, задержки, проблемы согласованности, наблюдаемости и распределённых операций.

\subsection{Сравнение}
\begin{longtable}{@{}p{0.23\textwidth}p{0.30\textwidth}p{0.37\textwidth}@{}}
\toprule
Архитектура & Идея & Особенность\\
\midrule
Серверный HTML & Сервер формирует HTML & Простая схема взаимодействия\\
Многослойное приложение & Разделение ответственности внутри сервера & Упрощает развитие кода\\
SPA + API & UI и API разделены & Один API может иметь разных клиентов\\
Монолит & Единое разворачиваемое приложение & Простота эксплуатации\\
Модульный монолит & Единое приложение с модулями & Структура без сетевой распределённости\\
Микросервисы & Несколько независимых сервисов & Гибкость ценой распределённой сложности\\
\bottomrule
\end{longtable}

\section{Часть II. Архитектура мастер-класса}
На занятии реализуется:
\begin{verbatim}
HTTP client
    |
    v
FastAPI
    |
    +-- Pydantic validation
    |
    +-- CRUD
    |
    v
SQLAlchemy ORM
    |
    v
SQLite
\end{verbatim}
Это небольшой монолитный REST-сервис. Его затем можно подключить к отдельному веб-клиенту.

\section{Часть III. HTTP, REST и CRUD}
Основные HTTP-методы:
\begin{center}
\begin{tabular}{ll}
\toprule
GET & получить ресурс\\
POST & создать ресурс\\
PUT & заменить ресурс\\
PATCH & частично изменить ресурс\\
DELETE & удалить ресурс\\
\bottomrule
\end{tabular}
\end{center}

Для ресурса \texttt{books}:
\begin{center}
\begin{tabular}{lll}
\toprule
Метод & URL & Назначение\\
\midrule
GET & /books & список книг\\
GET & /books/1 & книга с id=1\\
POST & /books & создание\\
PUT & /books/1 & изменение\\
DELETE & /books/1 & удаление\\
\bottomrule
\end{tabular}
\end{center}

\section{Часть IV. Мастер-класс}
\subsection{Предметная область}
Используем каталог книг. Поля: \texttt{id}, \texttt{title}, \texttt{author}, \texttt{year}.

\subsection{Шаг 1. Создание проекта}
\begin{lstlisting}[language=bash]
mkdir books-api
cd books-api
python -m venv .venv
\end{lstlisting}
Активация Windows PowerShell:
\begin{lstlisting}[language=bash]
.venv\Scripts\Activate.ps1
\end{lstlisting}
Linux/macOS:
\begin{lstlisting}[language=bash]
source .venv/bin/activate
\end{lstlisting}
Установка зависимостей:
\begin{lstlisting}[language=bash]
pip install fastapi uvicorn sqlalchemy
pip freeze > requirements.txt
\end{lstlisting}

\subsection{Шаг 2. Первый FastAPI endpoint}
Создать \texttt{main.py}:
\begin{lstlisting}[language=Python]
from fastapi import FastAPI

app = FastAPI(title="Books API")

@app.get("/")
def root():
    return {"message": "Books API"}
\end{lstlisting}
Запуск:
\begin{lstlisting}[language=bash]
uvicorn main:app --reload
\end{lstlisting}
Проверить \texttt{http://127.0.0.1:8000/} и документацию \texttt{/docs}.

\subsection{Шаг 3. Pydantic-модели}
Добавить:
\begin{lstlisting}[language=Python]
from pydantic import BaseModel

class BookCreate(BaseModel):
    title: str
    author: str
    year: int

class Book(BookCreate):
    id: int
\end{lstlisting}
\texttt{BookCreate} используется для входных данных, \texttt{Book} -- для ответа API.

\subsection{Шаг 4. GET списка: временное хранилище}
Для первого понимания CRUD можно сначала показать хранение в памяти:
\begin{lstlisting}[language=Python]
books: list[Book] = []
next_id = 1

@app.get("/books", response_model=list[Book])
def get_books():
    return books
\end{lstlisting}
После демонстрации сразу поставить вопрос: что произойдёт после перезапуска? Данные исчезнут. Значит, нужна БД.

\subsection{Шаг 5. SQLite и SQLAlchemy}
Добавить импорты:
\begin{lstlisting}[language=Python]
from sqlalchemy import create_engine, select
from sqlalchemy.orm import (
    DeclarativeBase,
    Mapped,
    Session,
    mapped_column,
    sessionmaker,
)
\end{lstlisting}
Настроить БД:
\begin{lstlisting}[language=Python]
DATABASE_URL = "sqlite:///./books.db"

engine = create_engine(
    DATABASE_URL,
    connect_args={"check_same_thread": False},
)

SessionLocal = sessionmaker(
    autocommit=False,
    autoflush=False,
    bind=engine,
)

class Base(DeclarativeBase):
    pass
\end{lstlisting}
ORM-модель:
\begin{lstlisting}[language=Python]
class BookDB(Base):
    __tablename__ = "books"

    id: Mapped[int] = mapped_column(
        primary_key=True,
        index=True,
    )
    title: Mapped[str]
    author: Mapped[str]
    year: Mapped[int]

Base.metadata.create_all(bind=engine)
\end{lstlisting}

\subsection{Шаг 6. Сессия БД}
\begin{lstlisting}[language=Python]
from collections.abc import Generator
from fastapi import Depends

def get_db() -> Generator[Session, None, None]:
    db = SessionLocal()
    try:
        yield db
    finally:
        db.close()
\end{lstlisting}
Смысл dependency: каждый запрос получает сессию, после завершения запроса сессия закрывается.

\subsection{Шаг 7. GET списка с пагинацией}
\begin{lstlisting}[language=Python]
from fastapi import Query

@app.get("/books", response_model=list[Book])
def get_books(
    skip: int = Query(0, ge=0),
    limit: int = Query(10, ge=1, le=100),
    db: Session = Depends(get_db),
):
    statement = (
        select(BookDB)
        .offset(skip)
        .limit(limit)
    )
    return db.scalars(statement).all()
\end{lstlisting}
Проверить:
\begin{verbatim}
GET /books
GET /books?skip=0&limit=2
GET /books?skip=2&limit=2
\end{verbatim}

\subsection{Шаг 8. GET одного объекта}
\begin{lstlisting}[language=Python]
from fastapi import HTTPException

@app.get("/books/{book_id}", response_model=Book)
def get_book(
    book_id: int,
    db: Session = Depends(get_db),
):
    book = db.get(BookDB, book_id)
    if book is None:
        raise HTTPException(
            status_code=404,
            detail="Book not found",
        )
    return book
\end{lstlisting}

\subsection{Шаг 9. POST}
\begin{lstlisting}[language=Python]
@app.post(
    "/books",
    response_model=Book,
    status_code=201,
)
def create_book(
    book_data: BookCreate,
    db: Session = Depends(get_db),
):
    book = BookDB(
        title=book_data.title,
        author=book_data.author,
        year=book_data.year,
    )
    db.add(book)
    db.commit()
    db.refresh(book)
    return book
\end{lstlisting}
Проверить через \texttt{/docs}, например:
\begin{lstlisting}
{
  "title": "Война и мир",
  "author": "Лев Толстой",
  "year": 1869
}
\end{lstlisting}

\subsection{Шаг 10. PUT}
\begin{lstlisting}[language=Python]
@app.put("/books/{book_id}", response_model=Book)
def update_book(
    book_id: int,
    book_data: BookCreate,
    db: Session = Depends(get_db),
):
    book = db.get(BookDB, book_id)
    if book is None:
        raise HTTPException(
            status_code=404,
            detail="Book not found",
        )

    book.title = book_data.title
    book.author = book_data.author
    book.year = book_data.year

    db.commit()
    db.refresh(book)
    return book
\end{lstlisting}

\subsection{Шаг 11. DELETE}
\begin{lstlisting}[language=Python]
@app.delete("/books/{book_id}", status_code=204)
def delete_book(
    book_id: int,
    db: Session = Depends(get_db),
):
    book = db.get(BookDB, book_id)
    if book is None:
        raise HTTPException(
            status_code=404,
            detail="Book not found",
        )

    db.delete(book)
    db.commit()
\end{lstlisting}

\section{Часть V. Полная финальная версия main.py}
Ниже находится версия, которую преподаватель может просто показать и выполнить без дописывания кода по ходу мастер-класса.

\begin{lstlisting}[language=Python]
from collections.abc import Generator

from fastapi import Depends, FastAPI, HTTPException, Query
from pydantic import BaseModel, Field
from sqlalchemy import create_engine, select
from sqlalchemy.orm import (
    DeclarativeBase,
    Mapped,
    Session,
    mapped_column,
    sessionmaker,
)

app = FastAPI(title="Books API")

DATABASE_URL = "sqlite:///./books.db"

engine = create_engine(
    DATABASE_URL,
    connect_args={"check_same_thread": False},
)

SessionLocal = sessionmaker(
    autocommit=False,
    autoflush=False,
    bind=engine,
)


class Base(DeclarativeBase):
    pass


class BookDB(Base):
    __tablename__ = "books"

    id: Mapped[int] = mapped_column(
        primary_key=True,
        index=True,
    )
    title: Mapped[str]
    author: Mapped[str]
    year: Mapped[int]


Base.metadata.create_all(bind=engine)


class BookCreate(BaseModel):
    title: str = Field(min_length=1, max_length=200)
    author: str = Field(min_length=1, max_length=200)
    year: int = Field(ge=0, le=2100)


class Book(BookCreate):
    id: int


def get_db() -> Generator[Session, None, None]:
    db = SessionLocal()
    try:
        yield db
    finally:
        db.close()


@app.get("/")
def root():
    return {"message": "Books API"}


@app.get("/books", response_model=list[Book])
def get_books(
    skip: int = Query(0, ge=0),
    limit: int = Query(10, ge=1, le=100),
    db: Session = Depends(get_db),
):
    statement = (
        select(BookDB)
        .offset(skip)
        .limit(limit)
    )
    return db.scalars(statement).all()


@app.get("/books/{book_id}", response_model=Book)
def get_book(
    book_id: int,
    db: Session = Depends(get_db),
):
    book = db.get(BookDB, book_id)
    if book is None:
        raise HTTPException(
            status_code=404,
            detail="Book not found",
        )
    return book


@app.post(
    "/books",
    response_model=Book,
    status_code=201,
)
def create_book(
    book_data: BookCreate,
    db: Session = Depends(get_db),
):
    book = BookDB(
        title=book_data.title,
        author=book_data.author,
        year=book_data.year,
    )
    db.add(book)
    db.commit()
    db.refresh(book)
    return book


@app.put("/books/{book_id}", response_model=Book)
def update_book(
    book_id: int,
    book_data: BookCreate,
    db: Session = Depends(get_db),
):
    book = db.get(BookDB, book_id)
    if book is None:
        raise HTTPException(
            status_code=404,
            detail="Book not found",
        )

    book.title = book_data.title
    book.author = book_data.author
    book.year = book_data.year

    db.commit()
    db.refresh(book)
    return book


@app.delete("/books/{book_id}", status_code=204)
def delete_book(
    book_id: int,
    db: Session = Depends(get_db),
):
    book = db.get(BookDB, book_id)
    if book is None:
        raise HTTPException(
            status_code=404,
            detail="Book not found",
        )

    db.delete(book)
    db.commit()
\end{lstlisting}

\section{Проверка мастер-класса}
Запуск:
\begin{lstlisting}[language=bash]
uvicorn main:app --reload
\end{lstlisting}
Открыть \texttt{http://127.0.0.1:8000/docs}.

Последовательность демонстрации:
\begin{enumerate}
    \item создать три книги;
    \item выполнить \texttt{GET /books};
    \item выполнить \texttt{GET /books/1};
    \item выполнить \texttt{GET /books/999} и показать 404;
    \item изменить книгу через \texttt{PUT};
    \item снова получить её через \texttt{GET};
    \item удалить её через \texttt{DELETE};
    \item снова получить её и показать 404;
    \item проверить пагинацию с разными \texttt{skip} и \texttt{limit};
    \item остановить сервер и запустить снова;
    \item показать, что данные сохранились в \texttt{books.db};
    \item отправить заведомо некорректные данные и показать автоматическую валидацию.
\end{enumerate}

Примеры данных:
\begin{lstlisting}
{
  "title": "Война и мир",
  "author": "Лев Толстой",
  "year": 1869
}
\end{lstlisting}
\begin{lstlisting}
{
  "title": "Преступление и наказание",
  "author": "Фёдор Достоевский",
  "year": 1866
}
\end{lstlisting}
\begin{lstlisting}
{
  "title": "Мастер и Маргарита",
  "author": "Михаил Булгаков",
  "year": 1967
}
\end{lstlisting}

\section{Педагогические акценты во время мастер-класса}
Не следует просто последовательно печатать код. Каждый существенный фрагмент вводится как ответ на проблему.

\begin{longtable}{@{}p{0.28\textwidth}p{0.62\textwidth}@{}}
\toprule
Проблема & Что вводится\\
\midrule
Нужно принимать HTTP-запросы & FastAPI\\
Нужно описать данные & Pydantic\\
Нужно хранить данные между запросами & SQLite\\
Нужно работать с БД из Python & SQLAlchemy ORM\\
Нужно получить конкретный объект & GET /books/\{id\}\\
Нужно создать объект & POST\\
Нужно изменить объект & PUT\\
Нужно удалить объект & DELETE\\
Нужно не возвращать тысячи строк & Пагинация\\
Нужно отклонять некорректные данные & Валидация\\
Нужно сообщать об отсутствии объекта & HTTP 404\\
\bottomrule
\end{longtable}

\section{Связь с практической работой №1}
К концу мастер-класса студенты видят полный минимальный путь:
\begin{verbatim}
HTTP request
    |
    v
FastAPI endpoint
    |
    v
Validation
    |
    v
CRUD operation
    |
    v
SQLAlchemy ORM
    |
    v
SQLite
    |
    v
HTTP response
\end{verbatim}

Практическая работа №1 требует реализовать ту же схему для другой предметной области.

\subsection{Минимальный чек-лист практической работы}
\begin{enumerate}
    \item выбрать вариант предметной области;
    \item определить сущности и их поля;
    \item создать ORM-модели;
    \item создать REST endpoints;
    \item реализовать CRUD;
    \item добавить валидацию входных данных;
    \item добавить обработку отсутствующих объектов;
    \item реализовать пагинацию/частичную загрузку;
    \item проверить API через Swagger UI или HTTP-клиент;
    \item убедиться, что данные сохраняются после перезапуска.
\end{enumerate}

\section{Вопросы студентам по ходу лекции}
\begin{enumerate}
    \item Чем серверный HTML отличается от REST API?
    \item Почему API можно использовать независимо от интерфейса?
    \item Может ли монолит быть хорошо спроектирован?
    \item Что нового появляется при переходе от монолита к микросервисам?
    \item Почему микросервисы не являются автоматически лучшим решением?
    \item Что произойдёт с данными, если хранить их в списке Python?
    \item Зачем нужна ORM?
    \item Зачем нужна серверная валидация, если клиент тоже может проверять данные?
    \item Зачем нужна пагинация?
    \item Почему запрос несуществующего объекта должен возвращать 404?
\end{enumerate}

\section{Итог первой лекции}
Главная мысль занятия:
\begin{quote}
Веб-приложение состоит не только из интерфейса. Клиент взаимодействует с сервером по HTTP; сервер принимает и валидирует запрос, выполняет операцию над предметными данными, обращается к базе данных и возвращает HTTP-ответ.
\end{quote}

При этом курс начинает с небольшого монолитного REST-сервиса. Далее система усложняется постепенно: сначала появляется отдельный веб-клиент, затем аутентификация и авторизация, после чего -- микросервисное и событийное взаимодействие.

\end{document}
